\begin{frame}[t]{\ev{\tagO}One hundred measurements can have the precision of nine}
\begin{columns}[T,onlytextwidth]
\column{.42\textwidth}
{\small\textbf{Calculation.} $N_{\rm eff}=N/(1+(N-1)\rho)$\par}
\vspace{.1cm}
{\small
\begin{tabular}{@{}rrr@{}}
\toprule
$\rho$ & $N=9$ & $N=100$\\
\midrule
0 & 9.0 & 100.0\\
0.05 & 6.4 & 16.8\\
0.10 & 5.0 & \textbf{9.2}\\
0.50 & 1.8 & 2.0\\
1.00 & 1.0 & 1.0\\
\bottomrule
\end{tabular}\par}
\column{.54\textwidth}
\fcolorbox{Line}{Pale}{\parbox{\dimexpr\linewidth-2\fboxsep-.8pt\relax}{\small\raggedright{\color{Teal}\bfseries SEPARATE EMPIRICAL COMPARISON}\\ Nine LLM judges from seven model families: Kish $N_{\rm eff}=2.18$, 95\% interval 2.07--2.31. A 2026 preprint on judge panels. It is not a fork measurement.}}\par
\vspace{.25cm}
{\small Kim et al.\ report another quantity: two models that both miss a question give the same wrong answer 60\% of the time, against 33\% by chance. That is a conditional agreement rate. It is not $\rho$.\par}
\end{columns}
\claim{At $\rho=0.1$, one hundred measurements carry the precision of about nine.}
\sources{Calculation from Kish (1965). \href{https://arxiv.org/abs/2605.29800}{Kohli, arXiv:2605.29800 (2026 preprint)}. \href{https://arxiv.org/abs/2506.07962}{Kim, Garg, Peng \& Garg, ICML 2025}, HELM, MMLU.}
\note{[0:45] One hundred measurements can have the precision of nine. At rho point one, one hundred correlated measurements give an effective size of nine point two. At rho one half, they give two. [hold up nine photocopies of one page] Nine copies of one page carry one page of information. Kohli's 2026 preprint finds nine judges worth about two point two independent votes. Kim and colleagues report a different quantity, a conditional agreement rate. The runtime can record where such sharing comes from.}
\end{frame}
